% Wave 16 S5/20 -- Seven faces of r_{CY} as (Sigma_2, C)-specialisations.
% Drinfeld + Etingof voice. Rectification C of two_stage_Phi_factorisation.
% Author: Raeez Lorgat. Subscripted kappa_*. No AI attribution.

\providecommand{\ClaimStatusTheorem}{\textsc{[Theorem]}}
\providecommand{\ClaimStatusConjectured}{\textsc{[Conjectured]}}
\providecommand{\ClaimStatusCorrected}{\textsc{[Corrected]}}
\providecommand{\ClaimStatusDefinition}{\textsc{[Definition]}}
\providecommand{\kBKM}{\kappa_{\mathrm{BKM}}}
\providecommand{\kch}{\kappa_{\mathrm{ch}}}
\providecommand{\kcat}{\kappa_{\mathrm{cat}}}
\providecommand{\kfib}{\kappa_{\mathrm{fiber}}}
\providecommand{\Sp}{\mathrm{Sp}}
\providecommand{\HolFA}{\mathrm{HolFA}}
\providecommand{\cF}{\mathcal{F}}
\providecommand{\cO}{\mathcal{O}}
\providecommand{\ZZ}{\mathbb{Z}}
\providecommand{\HH}{\mathbb{H}}
\providecommand{\PP}{\mathbb{P}}
\providecommand{\cA}{\mathcal{A}}
\providecommand{\rCY}{r_{\mathrm{CY}}}
\providecommand{\Muk}{\mathrm{Muk}}
\providecommand{\Niem}{\mathrm{Niem}}
\providecommand{\Hum}{\mathrm{Hum}}

\section{Seven faces of $\rCY$ as $(\Sigma_2,C)$-specialisations of
$\cF_{K3\times E}$}
\label{wn:sec:seven-faces-as-specialisations}

\emph{The seven arithmetic faces of $\rCY$ on $K3\times E$ are not seven
constructions. Five are $(\Sigma_2, C)$-specialisations of the single
$E_3$-hFA $\cF_{K3\times E}$; two are intrinsic invariants of the CY$_3$
that precede any specialisation. The multiplicity is a feature of the
specialisation datum, not a feature of $\Phi$.}

\subsection{The seven faces, enumerated}
\label{wn:subsec:seven-faces-enumerated}

For $X = K3\times E$, the arithmetic literature records seven
distinct numerical/algebraic incarnations of $\rCY$:
\begin{enumerate}
\item[(M)] $\rCY^{\Muk}$ --- Mukai-lattice companion of $\kcat$, reading
the Mukai vector pairing $\langle v_1, v_2\rangle = \chi(v_1, v_2)$ on
$\widetilde{H}(K3, \ZZ)$.
\item[(B)] $\rCY^{\mathrm{BKM}}$ --- Borcherds weight face: the
$c_1(0)/2 = 5$ weight of $\Delta_5$ (Gritsenko--Nikulin 1995, Borcherds
1998, Lorgat 2020).
\item[(F)] $\rCY^{\mathrm{fib}}$ --- K3-fibre rank $\kfib = 24$ (Mukai
lattice total rank of the K3 fibre).
\item[(H)] $\rCY^{\mathrm{Hodge}}$ --- $\kch(X) = \chi(\cO_X) = 0$
(Hodge supertrace face).
\item[(N)] $\rCY^{\Niem}$ --- Niemeier-lattice face: the 23 even
unimodular negative-definite rank-24 lattices indexing CHL quotients.
\item[(\Hum)] $\rCY^{\Hum}$ --- Humbert-divisor face:
monodromy of $\Delta_5$ around $H_1\cup H_4\subset\overline{\cA_2}$.
\item[(T)] $\rCY^{\mathrm{tw}}$ --- twined/CHL face: $(K3\times E)/(\ZZ/N\ZZ)$
quotients at $N\in\{1,2,3,4,6\}$ producing the paramodular Borcherds
products $\Phi_N$.
\end{enumerate}

\subsection{Three-tier hierarchy}
\label{wn:subsec:three-tier}

\begin{definition}[Three tiers]\ClaimStatusDefinition
\label{wn:def:three-tier}
The seven faces sort into three tiers by the layer of data they
probe:
\begin{center}
\renewcommand{\arraystretch}{1.22}
\begin{tabular}{|c|l|l|}
\hline
Tier & Layer & Faces \\
\hline
(i) & CY$_3$ intrinsic invariants (pre-specialisation) & (M), (H) \\
(ii) & Stage~1 invariants (the $E_3$-hFA $\cF_X$) & (F) \\
(iii) & $(\Sigma_2, C)$-specialisations of $\cF_X$ & (B), (N), (\Hum), (T) \\
\hline
\end{tabular}
\end{center}
\end{definition}

\paragraph{Tier (i): pre-specialisation CY invariants.}
(M) is the Mukai-lattice pairing on $K^\bullet(K3)$, a functorial invariant
of the $A_\infty$-category $D^b(K3)$ (Mukai 1988, \emph{Symplectic
structure on the moduli space of sheaves on an abelian or $K3$ surface},
Invent.\ Math.~77). (H) is $\chi(\cO_X) = \sum_q(-1)^q h^{0,q}$ on the
compact CY$_d$ total space, $0$ on $K3\times E$ by Künneth. Both read
the CY datum directly; neither consumes a $\Sigma_2$ or a reference
curve.

\paragraph{Tier (ii): the $E_3$-hFA itself.}
(F) is the fibre rank $\mathrm{rk}(H^{2,*}(K3))=24$ read off the $E_3$-hFA
$\cF_{K3\times E}\in E_3\text{-}\HolFA(K3\times E)$ before any $\Sigma_2$-integration.
It is the Hilbert-Chern character of the constant sheaf along the
K3 fibre direction of $\cF_X$ --- a property of $\cF_X$ as a Stage~1
object. No curve $C$ is needed: $\kfib$ is already a number.

\paragraph{Tier (iii): $(\Sigma_2, C)$-specialisations.}
Faces (B), (N), (\Hum), (T) distinguish themselves only once a
$\Sigma_2\subset X$ and a $C\subset X$ are fixed. They are
$\Sp_{\Sigma_2,C}(\cF_{X})$-invariants, not $\cF_X$-invariants.

\subsection{Tier (iii) face by face}
\label{wn:subsec:tier-iii-face-by-face}

\begin{conjecture}[BKM face as canonical specialisation]\ClaimStatusConjectured
\label{wn:conj:BKM-as-canonical-spec}
For $X = K3\times E$ and $(\Sigma_2, C) = (K3, E)$,
$$
\kBKM\bigl(\Sp_{K3, E}(\cF_{K3\times E})\bigr) = c_1(0)/2 = 5,
$$
where $c_1(0) = 10$ is the constant Fourier coefficient of $\phi_{0,1}$
(Lorgat 2020 \S2, Gritsenko--Nikulin alg-geom/9504006, Borcherds 1998
Thm.~13.3). This is the default $(\Sigma_2, C)$-choice: integrate $\cF$
over the K3 fibre, read the resulting $E_1$-chiral algebra on $E$,
identify its denominator automorphic form.
\end{conjecture}

\begin{conjecture}[Niemeier face as discrete $(\Sigma_2, C)$-twist]\ClaimStatusConjectured
\label{wn:conj:Niemeier-as-twist}
The 23 Niemeier lattices parametrise 23 distinct
$(\Sigma_2^{(i)}, C) = (K3^{(i)}, E)$-choices where the K3 fibre is
replaced by an isogenous K3 with Picard sublattice matching the
Niemeier-lattice root system. Each yields a distinct
$E_1$-chiral shadow on $E$ with Borcherds denominator form indexed by
the 23 automorphic-form orbits (Borcherds 1995, \emph{Automorphic forms
with singularities on Grassmannians}, Invent.~Math.~132). The 23
Niemeier lattices index $\Sigma_2$-choices, not $\Phi$-applications.
\end{conjecture}

\begin{conjecture}[Humbert face as $\Sigma_2$-monodromy]\ClaimStatusConjectured
\label{wn:conj:Humbert-as-monodromy}
The Humbert divisors $H_1, H_4\subset\overline{\cA_2}$ are the
boundary locus where the canonical $(\Sigma_2, C) = (K3, E)$-choice
degenerates: $H_1$ where $K3\to A_{\mathrm{KS}}$ factorises through
$E_1\times E_2$, $H_4$ through CM products. The Humbert-face invariant
is the local monodromy of $\Sp_{\Sigma_2, C}(\cF_X)$ as $(\Sigma_2, C)$
approaches the Humbert stratum: order~8 at $H_1$, order~2 at $H_4$
(Bruinier 2002 LNM~1780 Ch.~5; cf.\ \texttt{modular\_koszul\_bridge.tex}
Prop.~\texttt{prop:humbert-monodromy-K-kappa}). The face (\Hum) is the
\emph{discriminant locus} of the $(\Sigma_2, C)$-parameter space of
$\cF_X$.
\end{conjecture}

\begin{conjecture}[Twined face as CHL $(\Sigma_2, C)$-family]\ClaimStatusConjectured
\label{wn:conj:twined-as-CHL-family}
For each $N\in\{1,2,3,4,6\}$ the CHL quotient $X_N =
(K3\times E)/(\ZZ/N\ZZ)$ carries $\cF_{X_N}\in E_3\text{-}\HolFA(X_N)$.
The canonical $(\Sigma_2, C) = (K3/\langle g_N\rangle, E/\ZZ_N)$-specialisation
yields
$$
\kBKM\bigl(\Sp_{\Sigma_2^{(N)}, C^{(N)}}(\cF_{X_N})\bigr) = c_N(0)/2,
$$
with $c_N(0) \in \{10, 8, 6, 4, 2\}$ at $N\in\{1,2,3,4,6\}$ (Gritsenko
1999, \emph{Modular forms and moduli spaces of abelian and K3 surfaces};
Lorgat 2020 Conjecture~1 extends this to the 8 Gritsenko--Cléry forms).
The family of five CHL quotients is a $\ZZ_N$-equivariant family of
$(\Sigma_2, C)$-choices, not five separate functor invocations.
\end{conjecture}

\subsection{Cross-check: $\Sp_{\Sigma_2, C}(\cF_X)$ reproduces advertised structure}
\label{wn:subsec:cross-check}

\begin{conjecture}[Output matching, face by face]\ClaimStatusConjectured
\label{wn:conj:output-matching}
Under Rectification~C of Wave~15
(\texttt{two\_stage\_Phi\_factorisation.tex}
\S\texttt{wn:subsec:programme-rectification}), $\Sp_{\Sigma_2, C}(\cF_{K3\times E})$
outputs:
\begin{enumerate}
\item on $(K3, E)$: the Borcherds--Gritsenko $E_1$-chiral bialgebra
$\HH_{\Delta_5}$, matching face~(B).
\item on $(K3^{(i)}, E)$ Niemeier-twist: the Borcherds lift of the
Niemeier denominator form, matching face~(N).
\item on Humbert-limit of $(K3, E)$: the nearby-cycles $E_1$-chiral
carrying the order-8 monodromy, matching face~(\Hum).
\item on $(K3/\langle g_N\rangle, E/\ZZ_N)$: the Gritsenko paramodular
$E_1$-chiral, matching face~(T).
\end{enumerate}
Tier~(i) faces (M), (H) and tier~(ii) face (F) are not outputs of
$\Sp$; they are invariants of $X$ and $\cF_X$ upstream of
specialisation.
\end{conjecture}

\subsection{Consequence}
\label{wn:subsec:consequence-three-tier}

The seven faces are not seven siblings. Two are CY-datum intrinsics
(M, H); one is a Stage-1 invariant of $\cF_X$ (F); four are
$(\Sigma_2, C)$-specialisation data on a common Stage-1 object (B, N,
\Hum, T). The only genuine ``sibling'' relation is among the four
tier-(iii) faces, which are parametrised by a moduli space of
$(\Sigma_2, C)$-choices.

\begin{remark}[Face count is not a coincidence]
\label{wn:rmk:face-count}
The count ``seven'' is a numerology of the literature, not a
structural invariant. Under the three-tier decomposition: two CY
intrinsics, one $E_3$-hFA invariant, and the four-element moduli-space
stratum (canonical, Niemeier, Humbert-boundary, CHL-twisted) of
$(\Sigma_2, C)$-choices. The framework predicts further
tier-(iii) faces outside the canonical four (e.g.\ $\Sigma_2$ an
elliptic surface over $\PP^1$ rather than a K3; Wave~15~M3 frontier),
which would appear as Gritsenko--Cléry entries beyond the five-form
programme subset.
\end{remark}

\begin{remark}[Relation to the algebraic seven-face master]
\label{wn:rmk:part-VI-algebraic-faces}
The seven-face master chapter \ref{part:seven-faces-rCY},
file \texttt{cy\_holographic\_datum\_master.tex}, records
a \emph{different} seven-face enumeration: bar-cobar, CoHA, coisson,
MO $R$-matrix, affine super Yangian, elliptic Sklyanin, Gaudin.
Those are the \emph{algebraic presentations} of $\rCY$ as a binary
collision residue; the present enumeration is the \emph{arithmetic
witnesses} of the same $\rCY$ on $K3\times E$. Both sevens refer to the
same object; they are orthogonal slicings --- by algebraic
realisation versus by arithmetic specialisation. The three-tier
hierarchy refines only the arithmetic slicing.
\end{remark}
